% Problem record op_ad12295d8fdfb02a. This TeX file is derived from
% database/problems_json/op_ad12295d8fdfb02a.json by scripts/sync-tex.mjs; edit the JSON record
% and rerun the script rather than editing this file.
\section{Structural criterion for classical--entanglement trade-off advantage}
\label{sec:op_ad12295d8fdfb02a}

\paragraph{Problem.}
Characterize the finite-dimensional quantum channels for which joint
classical--entanglement coding strictly outperforms time sharing between
unassisted and unlimited-entanglement classical communication.  For a channel
$\mathcal N:A'\to B$, let $C_{\rm CE}(\mathcal N,e)$ be the supremum of
asymptotically achievable classical rates when at most $e$ ebits per channel
use are consumed.  Operationally,
\begin{equation}
  C_{\rm CE}(\mathcal N,e)
  :=\sup\left\{
    R:\begin{array}{l}
      \text{there are length-$n$ classical-message codes with}\\
      n^{-1}\log_2M_n\to R,\quad
      \limsup_{n\to\infty}n^{-1}\log_2K_n\leq e,\quad
      P_{\rm err}^{(n)}\to0
    \end{array}
  \right\},
  \label{eq:p79-ce-capacity}
\end{equation}
where $K_n$ in Eq.~\eqref{eq:p79-ce-capacity} is the Schmidt rank of the
preshared maximally entangled resource.  Define the unassisted and
unlimited-entanglement endpoints by
\begin{equation}
  C_0(\mathcal N):=C_{\rm CE}(\mathcal N,0),
  \qquad
  C_{\rm EA}(\mathcal N):=\sup_{e\geq0}C_{\rm CE}(\mathcal N,e),
  \qquad
  e_{\rm EA}(\mathcal N)
    :=\inf\{e:C_{\rm CE}(\mathcal N,e)=C_{\rm EA}(\mathcal N)\}.
  \label{eq:p79-ce-endpoints}
\end{equation}
For $e_{\rm EA}(\mathcal N)>0$, endpoint time sharing gives
\begin{equation}
  C_{\rm TS}(\mathcal N,e)
  :=\begin{cases}
    \left(1-\dfrac{e}{e_{\rm EA}(\mathcal N)}\right)C_0(\mathcal N)
    +\dfrac{e}{e_{\rm EA}(\mathcal N)}C_{\rm EA}(\mathcal N),
      &0\leq e\leq e_{\rm EA}(\mathcal N),\\[2mm]
    C_{\rm EA}(\mathcal N),&e\geq e_{\rm EA}(\mathcal N).
  \end{cases}
  \label{eq:p79-time-sharing-curve}
\end{equation}
If the infimum in Eq.~\eqref{eq:p79-ce-endpoints} is not attained, interpret
Eq.~\eqref{eq:p79-time-sharing-curve} through its operational closure; if
$e_{\rm EA}(\mathcal N)=0$, set
$C_{\rm TS}(\mathcal N,e)=C_{\rm EA}(\mathcal N)$.  Give intrinsic
necessary and sufficient conditions for strict suboptimality of this
benchmark:
\begin{equation}
  \exists e>0:\qquad
  C_{\rm CE}(\mathcal N,e)>C_{\rm TS}(\mathcal N,e).
  \label{eq:p79-strict-ce-advantage}
\end{equation}

\subsection*{Status}

Unsolved

\subsection*{Source}

Wilde explicitly asks how to determine which channels benefit from trade-off
coding rather than time sharing in Section~22.5 of his text
\sourcecite{ref:p79-wilde}{Wil17}.  Br\'adler, Hayden, Touchette, and Wilde
likewise call for a general method that determines the gain over time sharing
\sourcecite{ref:p79-bradler-et-al}{BHTW10}.

\subsection*{Progress}

\begin{itemize}
  \item Hsieh and Wilde proved a matching achievable region and multiletter
  converse for simultaneous classical communication, quantum communication,
  and entanglement.  The classical--entanglement slice determines
  $C_{\rm CE}(\mathcal N,e)$ through a regularized optimization, but it does
  not provide an intrinsic channel-level criterion for
  Eq.~\eqref{eq:p79-strict-ce-advantage}
  \sourcecite{ref:p79-hsieh-wilde}{HW10}.

  \item Br\'adler, Hayden, Touchette, and Wilde derived exact trade-off
  regions for Hadamard channels.  They exhibit strict improvement over
  endpoint time sharing in nontrivial parameter regimes, including
  finite-dimensional dephasing and $1\to N$ cloning channels, and introduce
  a quantitative measure of the gain.  The trivial parameter endpoints need
  not show strict advantage, and the examples do not characterize all
  finite-dimensional channels
  \sourcecite{ref:p79-bradler-et-al}{BHTW10}.
\end{itemize}

\subsection*{References}

\begin{enumerate}[label={},leftmargin=4.8em,labelsep=0.5em]
  \footnotesize
  \item[\textup{[Wil17]}]\label{ref:p79-wilde}
  M. M. Wilde, \emph{Quantum Information Theory}, 2nd ed., Cambridge
  University Press (2017), Section~22.5.
  \href{https://doi.org/10.1017/9781316809976}{doi:10.1017/9781316809976};
  \href{https://arxiv.org/abs/1106.1445}{arXiv:1106.1445}.

  \item[\textup{[HW10]}]\label{ref:p79-hsieh-wilde}
  M.-H. Hsieh and M. M. Wilde,
  ``Trading Classical Communication, Quantum Communication, and Entanglement
  in Quantum Shannon Theory,''
  \emph{IEEE Trans. Inf. Theory} \textbf{56}, 4705--4730 (2010).
  \href{https://doi.org/10.1109/TIT.2010.2054532}{doi:10.1109/TIT.2010.2054532};
  \href{https://arxiv.org/abs/0901.3038}{arXiv:0901.3038}.

  \item[\textup{[BHTW10]}]\label{ref:p79-bradler-et-al}
  K. Br\'adler, P. Hayden, D. Touchette, and M. M. Wilde,
  ``Trade-Off Capacities of the Quantum Hadamard Channels,''
  \emph{Physical Review A} \textbf{81},
  062312 (2010).
  \href{https://doi.org/10.1103/PhysRevA.81.062312}{doi:10.1103/PhysRevA.81.062312};
  \href{https://arxiv.org/abs/1001.1732}{arXiv:1001.1732}.
\end{enumerate}

\subsection*{Comment}

Strict trade-off advantage itself is known to occur.  The unresolved problem
is a necessary-and-sufficient structural characterization, including any
regularization across tensor powers, that identifies when endpoint time
sharing is optimal.  The statement is deliberately
restricted to the precisely defined classical--entanglement trade-off; it
does not use an informal full-dynamic time-sharing region.

\subsection*{Field}

Quantum Communication; Quantum Resource Theory

\subsection*{Topic}

Classical capacity; Entanglement-assisted communication; Additivity and regularization

\subsection*{ID}

\texttt{op\_ad12295d8fdfb02a}
